% Generated. Do not edit manually. Edit resume_content.json instead.
\renewcommand{\ResumeSummary}{
Ph.D. mechanical engineer combining nonlinear finite element methods, multiphysics simulation, and experimental model calibration. Developed custom FEM/XFEM solvers, nonlinear material models, and a neural-network damage surrogate trained on simulation data. Applied structural and impact analysis to product development, built Python-based engineering automation tools, and validated designs through prototype testing, with additional experience developing real-time C/C++ simulation software.
}

\renewcommand{\ResumeSkills}{
\textbf{Computational Mechanics:} FEM/XFEM, Nonlinear material modeling, Explicit dynamic FEA\par
\textbf{Simulation Methods:} Neural-network damage surrogates, Reduced-order modeling, Numerical solver development, Structural optimization\par
\textbf{Engineering Automation:} Python, CadQuery, C/C++, Parametric CAD, Jenkins CI/CD\par
\textbf{Experimental Validation:} Prototype testing, Impact-hammer testing, Structural model calibration\par
}

\renewcommand{\RheemBullets}{
\item \textbf{Explicit Dynamic Impact and Fracture Assessment}
\begin{itemize}[leftmargin=1.3em,topsep=1pt,itemsep=1pt]
\item Performed explicit dynamic FEA impact analysis of plastic components to evaluate impact-induced stresses and potential cracking; calculated the J-integral to assess the driving force for crack propagation.
\end{itemize}
\item \textbf{Airflow-Induced Condensate Management Invention}
\begin{itemize}[leftmargin=1.3em,topsep=1pt,itemsep=1pt]
\item Invented a passive condensate entrainment and slinging system that uses fan-generated pressure differences to draw condensed water into the fan airflow path and improve condensate slinging performance.
\item Performed CFD simulations to investigate airflow behavior, pressure distribution, and condensate entrainment mechanisms.
\item Designed, fabricated, and iteratively refined 3D-printed prototypes to demonstrate the mechanism.
\item Conducted experimental testing to validate condensate entrainment, slinging effectiveness, and overall system performance.
\end{itemize}
\item \textbf{Motor-Mount Structural Optimization}
\begin{itemize}[leftmargin=1.3em,topsep=1pt,itemsep=1pt]
\item Redesigned a motor-mount structure using FEA, reducing material usage by more than 50\% while maintaining stiffness and vibration requirements. Simplified injection-molded components to reduce secondary manufacturing operations and assembly complexity.
\end{itemize}
\item \textbf{Structural FEA and Assembly Optimization of Modular AHU}
\begin{itemize}[leftmargin=1.3em,topsep=1pt,itemsep=1pt]
\item Performed comprehensive finite element analysis (FEA) of the nine-segment modular AHU assembly to evaluate structural integrity, press-fit connections, and load-bearing capacity. Analyzed deformation under fan and motor loads to maintain critical clearances between the rotating fan lip and stationary bellmouth, preventing mechanical interference and supporting reliable prototype operation.
\end{itemize}
\item \textbf{Multi-Material Product FEA Analysis}
\begin{itemize}[leftmargin=1.3em,topsep=1pt,itemsep=1pt]
\item Performed finite element analysis (FEA) to support product designs using ABS and ASA plastics, aluminum and steel sheet metal, and expanded polypropylene (EPP) at different densities.
\end{itemize}
\item \textbf{Parametric Fan-Shroud Generation}
\begin{itemize}[leftmargin=1.3em,topsep=1pt,itemsep=1pt]
\item Developed a Python/CadQuery-based tool that converts fan diameter, coil geometry, component locations, and design clearances into 3D shroud geometry and STEP files, reducing repetitive CAD modeling and accelerating design iteration across product configurations.
\end{itemize}
}

\renewcommand{\LTTSBullets}{
\item \textbf{Modular GUI for Aerodynamic Data Analysis}
\begin{itemize}[leftmargin=1.3em,topsep=1pt,itemsep=1pt]
\item Developed a modular GUI for aerodynamic data analysis with well-defined input/output interfaces, improving maintainability and reducing computational and memory overhead.
\end{itemize}
\item \textbf{Real-Time C/C++ HVAC Plant Model for SiL and HiL}
\begin{itemize}[leftmargin=1.3em,topsep=1pt,itemsep=1pt]
\item Developed a real-time C/C++ executable HVAC plant model for Software-in-the-Loop (SiL) and Hardware-in-the-Loop (HiL) applications, combining finite element solution of transient heat transfer and velocity-pressure-coupled Navier-Stokes equations with lumped-capacitance heater/cooler models, convective thermal-fluid coupling, and a proportional-derivative (PD) HVAC controller model.
\end{itemize}
\item \textbf{Jenkins CI/CD and Automated Debugging}
\begin{itemize}[leftmargin=1.3em,topsep=1pt,itemsep=1pt]
\item Developed Jenkins CI/CD pipelines with Bash and PowerShell scripting and automated screen-capture capability, streamlining debugging and test workflows for construction and forestry vehicle software.
\end{itemize}
\item \textbf{Probability-Based File-Type Detection Algorithm}
\begin{itemize}[leftmargin=1.3em,topsep=1pt,itemsep=1pt]
\item Developed a probability-based file-type detection algorithm for John Deere software, replacing a rigid if-else structure with a scalable, data-driven approach intended to improve performance, reliability, and extensibility.
\end{itemize}
}

\renewcommand{\UTSABullets}{
\item \textbf{XFEM-Informed Neural Network Damage Surrogate}
\begin{itemize}[leftmargin=1.3em,topsep=1pt,itemsep=1pt]
\item Developed a neural-network surrogate trained on high-fidelity XFEM simulations to approximate localized damage behavior in larger structural models.
\end{itemize}
\item \textbf{XFEM, Nonlinear Materials, and Damage Modeling}
\begin{itemize}[leftmargin=1.3em,topsep=1pt,itemsep=1pt]
\item Developed XFEM methods to represent stress fields around three-dimensional voids and implemented nonlinear material behavior incorporating plasticity, strain-rate effects, thermal expansion, and von Mises and Drucker-Prager yield criteria.
\end{itemize}
\item \textbf{Experimental Structural Digital Twin and Model Updating}
\begin{itemize}[leftmargin=1.3em,topsep=1pt,itemsep=1pt]
\item Created a structural digital twin of an experimental habitat by acquiring impact-hammer and accelerometer data, identifying natural frequencies and mode shapes, using Crypton optical geometry measurements, and calibrating a parameterized beam-element model to match identified modal properties and measured acceleration histories.
\end{itemize}
\item \textbf{Custom 3D Multiphysics Finite Element Solvers}
\begin{itemize}[leftmargin=1.3em,topsep=1pt,itemsep=1pt]
\item Developed custom three-dimensional finite element solvers from first principles for structural mechanics, transient heat transfer, and fluid dynamics, including coupled structural-thermal habitat models and interior thermal-fluid simulations.
\end{itemize}
\item \textbf{Reduced-Order Modeling and High-Order Elements}
\begin{itemize}[leftmargin=1.3em,topsep=1pt,itemsep=1pt]
\item Developed modal-coordinate reduced-order models and applied Guyan static condensation; implemented very high-order finite element formulations for acoustic wave propagation associated with micrometeoroid impacts.
\end{itemize}
\item \textbf{Efficient Scientific Software and Numerical Algorithms}
\begin{itemize}[leftmargin=1.3em,topsep=1pt,itemsep=1pt]
\item Developed computational tools emphasizing reduced arithmetic and memory complexity to support novel numerical formulations and simulation workflows; contributed to a publication on algorithm development.
\end{itemize}
\item \textbf{NASA Modular-Coupled Virtual Testbed (MCVT)}
\begin{itemize}[leftmargin=1.3em,topsep=1pt,itemsep=1pt]
\item Developed and integrated structural, thermal, and fluid plant models for NASA's Modular-Coupled Virtual Testbed (MCVT), a system-of-systems simulation environment for virtual testing of resilient extraterrestrial habitats.
\end{itemize}
\item \textbf{Thermal Protective Layer Topology Optimization}
\begin{itemize}[leftmargin=1.3em,topsep=1pt,itemsep=1pt]
\item Investigated thermo-mechanical topology optimization of lunar habitat thermal protective layers under different environmental conditions and lunar latitudes.
\end{itemize}
}

\renewcommand{\OtherExperience}{

}

\renewcommand{\Demonstrations}{
Large-deformation beam simulator: \url{https://www.linkedin.com/feed/update/urn:li:activity:7429372567318331392/}\par
Finite element formulation software tool: \url{https://www.linkedin.com/feed/update/urn:li:activity:7122914707790843904/}\par
MCVT structural-mechanical model: \url{https://youtu.be/TAA66P_UXIM}\par
}

\let\ResumeDemonstrations\Demonstrations\renewcommand{\Demonstrations}{}\renewcommand{\OtherExperience}{}
